\section{Positive packets and uniform relative estimates}\label{sec:positivity}

Throughout this section,
\begin{equation}\label{eq:time-range}
 0\le t\le\frac12,\qquad x\in\R.
\end{equation}
The restriction on time is used explicitly in the convexity proof.
Complex spectral arguments will be written $z=x+iy$ with $x,y\in\R$.
The exact arithmetic identities of Section~\ref{sec:packets} themselves
hold for every real heat time.

\subsection{Strict convexity of the higher source modes}

\begin{theorem}[Positivity of the arithmetic modes]\label{thm:positive}
For every $n\ge2$, $H_n(t,x)>0$ on \eqref{eq:time-range}.
More generally, if $(c_n)_{n\ge2}$ is bounded, nonnegative and not identically
zero, then
\[
 \int_0^\infty e^{tu^2}\Bigl(\sum_{n\ge2}c_n\phi_n(u)\Bigr)
                  \cos(xu)\dd u>0.
\]
\end{theorem}
\begin{proof}
Put $g(u)=e^{tu^2}\phi_n(u)$ and $a=\pi n^2e^{4u}$. For $n\ge2$ and
$u\ge0$, $a>12$ and $g>0$. Since
$g=e^{tu^2+u-a}a(2a-3)$, logarithmic differentiation gives
\begin{align}
 \ell'(u)&=2tu+1-4a+4+\frac{8a}{2a-3}
          =2tu+9-4a+\frac{12}{2a-3},\label{eq:log-first}\\
 \ell''(u)&=2t-16a-\frac{96a}{(2a-3)^2},
 \qquad \ell=\log g.\label{eq:log-second}
\end{align}
The elementary inequalities
\[
 2tu\le u,\qquad \frac{12}{2a-3}<\frac47,\qquad
 a\ge12(1+4u)>u+9+\frac47
\]
show that $\ell'\le-3a$. Also
$96/(2a-3)^2<96/441<1$, so $\ell''\ge-17a$. Therefore
\begin{equation}\label{eq:convexity}
 g'<0,\qquad
 \frac{g''}{g}=(\ell')^2+\ell''
 \ge9a^2-17a\ge7a^2>0.
\end{equation}
The last inequality uses $2a\ge17$, which follows from $a>12$.

All boundary terms at infinity vanish by Lemma~\ref{lem:regularity}.
For $x\ne0$, two integrations by parts give
\[
 \int_0^\infty g(u)\cos(xu)\dd u
 =-\frac{g'(0)}{x^2}
  -\frac1{x^2}\int_0^\infty g''(u)\cos(xu)\dd u.
\]
Because $\int_0^\infty g''(u)\dd u=-g'(0)$, this is
\begin{equation}\label{eq:positive-integral}
 \boxed{\int_0^\infty g(u)\cos(xu)\dd u
 =\frac1{x^2}\int_0^\infty g''(u)(1-\cos(xu))\dd u>0.}
\end{equation}
Strict positivity follows from $g''>0$ and from the fact that
$1-\cos(xu)$ is positive on some interval. At $x=0$, positivity follows
from the defining integral. For bounded nonnegative coefficients, absolute
convergence permits summation of the integrals; at least one summand is
strictly positive.
\end{proof}

Formula \eqref{eq:positive-integral} is the elementary convex-kernel
principle underlying P\'olya's criterion; compare
\cite[Theorem 3.4]{Csordas}. Its application here depends on the explicit
source polynomial in \eqref{eq:source}. The calculation did not establish
convexity for $n=1$ or for the complete source.

\begin{corollary}\label{cor:positive-packets}
If $\mathcal P$ is a nonempty finite prime set, then
$A_{\mathcal P}$, $R_{\mathcal P}$ and $C$ are strictly positive on the
real axis throughout \eqref{eq:time-range}. Also $G_d>0$ for every $d\ge2$.
\end{corollary}
\begin{proof}
Each is a nonempty sum of terms $H_n$ with $n\ge2$. The remainder is
nonempty because $Q+1>1$ is relatively prime to $Q$. The theorem applies
to each indicator coefficient sequence.
\end{proof}

In particular, the minus sign in $G_2+G_3-G_6$ is harmless for positivity
only because the complete expression is the coefficient selector in
\eqref{eq:A}. This does not assign signs to its separated complex-shift
pieces.

\subsection{Positive real part in a complex strip}

For an integer $m\ge2$ and a bounded nonnegative coefficient sequence
$(c_n)_{n\ge m}$ which is not identically zero, write
\begin{equation}\label{eq:positive-combination}
 V(t,z)=\sum_{n\ge m}c_nH_n(t,z).
\end{equation}
Lemma~\ref{lem:regularity} shows that this is entire in $z$, with
termwise differentiation valid on compact sets.

\begin{theorem}[A positive-real-part strip]\label{thm:strip-positive}
For every combination \eqref{eq:positive-combination} and every time in
\eqref{eq:time-range},
\begin{equation}\label{eq:positive-strip}
 \Re V(t,x+iy)>0\qquad(x\in\R,\ |y|\le\pi m^2).
\end{equation}
In particular, $V$ has no zero in this closed strip.
\end{theorem}
\begin{proof}
For each $n\ge m$, put $g=e^{tu^2}\phi_n$ and
$a=\pi n^2e^{4u}$. Equations \eqref{eq:log-first} and
\eqref{eq:convexity} give
\[
 -4a<\frac{g'}g\le-3a,\qquad \frac{g''}g\ge7a^2.
\]
For fixed real $y$ define $k_y(u)=g(u)\cosh(yu)$. Since
$0\le y\tanh(yu)\le|y|$ for $u\ge0$,
\begin{align}
 \frac{k_y'}{k_y}&\le-3a+|y|<0,\notag\\
 \frac{k_y''}{k_y}
 &=\frac{g''}g+2\frac{g'}g\,y\tanh(yu)+y^2\notag\\
 &\ge7a^2-8a|y|+y^2
   =(a-|y|)(7a-|y|).\label{eq:strip-convexity}
\end{align}
Here $|y|\le\pi m^2\le\pi n^2\le a$. The last expression is
strictly positive for $u>0$, including when $|y|=\pi m^2$.
All required boundary terms at infinity vanish by
Lemma~\ref{lem:regularity}. Because
$\Re\cos((x+iy)u)=\cos(xu)\cosh(yu)$, the same two integrations by
parts as in \eqref{eq:positive-integral} yield, for $x\ne0$,
\[
 \Re H_n(t,x+iy)
 =\frac1{x^2}\int_0^\infty k_y''(u)(1-\cos(xu))\dd u>0.
\]
At $x=0$, the defining integral is positive. Bounded coefficients and
the decay estimate justify summing the integrals, and at least one
coefficient is positive. This proves \eqref{eq:positive-strip}.
\end{proof}

Thus $A_{\mathcal P}$ and $C$ have positive real part for $|y|\le4\pi$,
and a remainder starting at $m$ has positive real part for
$|y|\le\pi m^2$. This property uses the positive higher-mode
coefficients. It imposes neither primality nor full source reflection.

\subsection{A quantitative lower bound}

Define the fixed positive constant
\begin{equation}\label{eq:kappa}
 \kappa=\frac78\phi_2(0)
       =\frac78(32\pi^2-12\pi)e^{-4\pi}.
\end{equation}

\begin{lemma}[A sinc estimate]\label{lem:sinc}
With the continuous value at $v=0$,
\begin{equation}\label{eq:sinc}
 1-\frac{\sin v}{v}\ge\frac{v^2}{6(1+v^2)}
 \qquad(v\in\R).
\end{equation}
\end{lemma}
\begin{proof}
Both sides are even. For $0\le v\le2$, the alternating Taylor estimate
gives $1-\sin v/v\ge v^2/6-v^4/120$. This is at least the right side
of \eqref{eq:sinc}, since
$(20-v^2)(1+v^2)-20=v^2(19-v^2)\ge0$.
For $v\ge2$, $1-\sin v/v\ge1-1/v\ge1/2$, whereas the right side is
at most $1/6$.
\end{proof}

\begin{proposition}[A nonvanishing arithmetic carrier]\label{prop:lower}
If $2\in\mathcal P$, then
\begin{equation}\label{eq:lower}
 A_{\mathcal P}(t,x)\ge H_2(t,x)
       \ge\frac{\kappa}{4096+x^2}.
\end{equation}
\end{proposition}
\begin{proof}
Only the second inequality needs proof. Let $L=1/64$, $n=2$ and
$g=e^{tu^2}\phi_2$. For $0\le u\le L$,
\[
 a=4\pi e^{4u}\le4\pi e^{1/16}
 <\frac{88}{7}\frac{16}{15}<14.
\]
Here $e^{1/16}<16/15$ follows by comparing its power series with the
geometric series, and $\pi<22/7$. Equation \eqref{eq:log-first} implies
$\ell'>9-4a>-47$. Hence
\[
 g(u)\ge\phi_2(0)e^{-47/64}>\frac{\phi_2(0)}3,
\]
where we used $47/64<1$ and $e<3$. Equation \eqref{eq:convexity} now gives
$g''(u)\ge7\cdot12^2\phi_2(0)/3=336\phi_2(0)$ on this interval.
For $x\ne0$, retain only this interval in \eqref{eq:positive-integral}:
\begin{align*}
 H_2(t,x)
 &\ge\frac{336\phi_2(0)}{x^2}
       \left(L-\frac{\sin(xL)}x\right)\\
 &\ge\frac{56\phi_2(0)L^3}{1+x^2L^2}
  =\frac{(7/8)\phi_2(0)}{4096+x^2}.
\end{align*}
The same inequality holds at $x=0$ by continuity. Since $n=2$ belongs to
the selector defining $A_{\mathcal P}$, all its other terms are positive.
\end{proof}

The denominator in \eqref{eq:lower} is deliberate: the packets have an
inverse-square tail, unlike the rapidly decreasing complete transform.
This estimate permits normalization without dividing by $H$ or $H_x$.

The corresponding quantitative strip estimate uses the constant
\begin{equation}\label{eq:kappa-strip}
 \kappa_s=\frac34\phi_2(0)=\frac67\kappa.
\end{equation}

\begin{proposition}[A lower bound in the unit strip]\label{prop:strip-lower}
If $2\in\mathcal P$, then, for $0\le t\le1/2$,
\begin{equation}\label{eq:strip-lower}
 \Re A_{\mathcal P}(t,x+iy)
 \ge\Re H_2(t,x+iy)
 \ge\frac{\kappa_s}{4096+x^2}
 \qquad(x\in\R,\ |y|\le1).
\end{equation}
\end{proposition}
\begin{proof}
Use $g=e^{tu^2}\phi_2$, $k_y=g\cosh(yu)$ and $L=1/64$.
For $|y|\le1$ and $a\ge12$, \eqref{eq:strip-convexity} gives
\[
 \frac{k_y''}{k_y}\ge7a^2-8a\ge6a^2.
\]
The local estimate in Proposition~\ref{prop:lower} is
$g(u)\ge\phi_2(0)/3$ on $0\le u\le L$. Since $\cosh(yu)\ge1$,
it follows that $k_y''(u)\ge288\phi_2(0)$ there. For $x\ne0$,
the positive-integral formula and Lemma~\ref{lem:sinc} give
\[
 \Re H_2(t,x+iy)
 \ge\frac{288\phi_2(0)}{x^2}
       \left(L-\frac{\sin(xL)}x\right)
 \ge\frac{48\phi_2(0)L^3}{1+x^2L^2}
 =\frac{(3/4)\phi_2(0)}{4096+x^2}.
\]
Continuity gives the same estimate at $x=0$. Every other term in
$A_{\mathcal P}$ has positive real part by
Theorem~\ref{thm:strip-positive}, which proves the first inequality.
\end{proof}

\subsection{Endpoint masses and positive tail bounds}

For an integer $m\ge2$, use $r_{m,k}$ from \eqref{eq:strip-tail} and set
\begin{align}
 S_m&=\frac{2\pi^2m^4e^{-\pi m^2}}{1-r_{m,4}},&
 Q_m&=\frac{8\pi^3m^6e^{-\pi m^2}}{1-r_{m,6}},\label{eq:SQ}\\
 \epsilon_j(m)&=\frac{S_m j!}{(4\pi m^2-9)^{j+1}}
 \quad(j\ge0).\label{eq:epsilon}
\end{align}
The denominators are positive: for example,
$r_{m,6}\le(3/2)^6e^{-5\pi}<1$ when $m\ge2$.
Let $S$ be a nonempty subset of $\{m,m+1,\ldots\}$ and write
\[
 g_S(u)=e^{tu^2}\sum_{n\in S}\phi_n(u),\qquad
 V(t,x)=\int_0^\infty g_S(u)\cos(xu)\dd u.
 \]
The use of $S$ as an index set is distinct from the scalar bound $S_m$.

\begin{lemma}\label{lem:mass}
With $q_S=-g_S'(0)$ and $s_S=g_S(0)$, one has
\begin{align}
 0<s_S&\le S_m,&0<q_S&\le Q_m,&
 \int_0^\infty u^j g_S(u)\dd u&\le\epsilon_j(m),\label{eq:mass-bounds}\\
 \int_0^\infty g_S''\dd u&=q_S,&
 \int_0^\infty u g_S''\dd u&=s_S,&
 \int_0^\infty u^2g_S''\dd u&=2\int_0^\infty g_S\dd u.\label{eq:mass-identities}
\end{align}
Moreover
\begin{equation}\label{eq:two-upper}
 0<V(t,x)\le\epsilon_0(m),\qquad
 0<V(t,x)\le\frac{2Q_m}{x^2}\quad(x\ne0).
\end{equation}
\end{lemma}
\begin{proof}
The moment bound is Lemma~\ref{lem:strip-tail} at $T_+=1/2$, $Y=0$,
or directly its nonoscillatory majorant. At $u=0$,
$\phi_n(0)\le2\lambda_n^2e^{-\lambda_n}$ and
\[
 -\phi_n'(0)
 =(8\lambda_n^3-30\lambda_n^2+15\lambda_n)e^{-\lambda_n}
 <8\lambda_n^3e^{-\lambda_n}.
\]
The derivative of $e^{tu^2}$ is zero at $0$, so it does not change this
endpoint mass. Summing successive-ratio bounds gives $S_m$ and $Q_m$.
Positivity follows from $g_S>0$ and $g_S'<0$. Integration by parts gives
the three identities in \eqref{eq:mass-identities}; all upper endpoints
vanish. For the last identity, first integrate $u^2g_S''$ to obtain
$-2\int u g_S'$, and integrate once more. Finally,
\eqref{eq:positive-integral}, applied to the summed source, and
$0\le1-\cos(xu)\le2$ give the inverse-square upper bound.
The absolute source integral gives the other one.
\end{proof}

\begin{theorem}[Uniform relative approximation]\label{thm:relative}
Suppose $2\in\mathcal P$ and every index in $R_{\mathcal P}$ is at least
$m\ge2$. Define
\begin{equation}\label{eq:rho}
 \rho_m=\frac{4096\epsilon_0(m)+2Q_m}{\kappa}.
\end{equation}
Then, throughout \eqref{eq:time-range},
\begin{equation}\label{eq:relative}
 0<r_{\mathcal P}:=\frac{R_{\mathcal P}}{A_{\mathcal P}}\le\rho_m,
 \qquad
 0<\frac{C-A_{\mathcal P}}{C}
      \le\frac{\rho_m}{1+\rho_m}<\rho_m.
\end{equation}
In particular,
\[
 \rho_3<0.000263,\qquad \rho_5<7.3\times10^{-25},
\]
applicable respectively to $\mathcal P=\{2\}$ and $\mathcal P=\{2,3\}$.
\end{theorem}
\begin{proof}
Equations \eqref{eq:lower} and \eqref{eq:two-upper} give
\[
 \kappa\frac{R_{\mathcal P}}{A_{\mathcal P}}
 \le(4096+x^2)R_{\mathcal P}
 \le4096\epsilon_0(m)+2Q_m.
\]
For $x=0$ the same bound follows from $R_{\mathcal P}\le\epsilon_0(m)$.
Since $C=A_{\mathcal P}(1+r_{\mathcal P})$, the second bound follows by
monotonicity of $r/(1+r)$ for $r\ge0$. The numerical inequalities are
verified from the explicit constants in
Section~\ref{sec:certification}, without a height cutoff.
\end{proof}

\subsection{Logarithmic derivatives and remainder jets}

The positive-real-part strip controls two real logarithmic derivatives
at once. Retaining the relation between these two derivatives gives a
stronger quotient estimate than bounding them separately.

\begin{lemma}[A Schwarz coefficient bound]\label{lem:strip-Schwarz}
For $V$ as in \eqref{eq:positive-combination}, define
$\beta(t,x)=V_x(t,x)/V(t,x)$. Then
\begin{equation}\label{eq:strip-log-budget}
 \beta(t,x)^2+|\beta_x(t,x)|\le\frac1{4m^4}
 \qquad(0\le t\le1/2,\ x\in\R).
\end{equation}
In particular, $|\beta|\le1/(2m^2)$.
\end{lemma}
\begin{proof}
Fix $t,x$ and put $L=\pi m^2$, $K=4L/\pi=4m^2$. The map
$\zeta\mapsto x+K\operatorname{artanh}\zeta$ takes the unit disk
into $|\Im z|<L$: indeed
$\operatorname{artanh}\zeta=\frac12\log((1+\zeta)/(1-\zeta))$
has imaginary part between $-\pi/4$ and $\pi/4$. Define
\[
 f(\zeta)=\frac{V(t,x+K\operatorname{artanh}\zeta)}{V(t,x)},\qquad
 q(\zeta)=\frac{f(\zeta)-1}{f(\zeta)+1}.
\]
Theorem~\ref{thm:strip-positive} and $V(t,x)>0$ imply
$\Re f>0$, $|q|<1$ and $q(0)=0$. Write
$q(\zeta)=a_1\zeta+a_2\zeta^2+O(\zeta^3)$. The coefficients are
real by conjugation symmetry. Schwarz's lemma makes
$h(\zeta)=q(\zeta)/\zeta$, with its removable value at zero, a map
into the closed unit disk. If $|a_1|<1$, apply Schwarz's lemma to
$(h-a_1)/(1-a_1h)$ to obtain
\[
 |a_2|\le1-a_1^2.
\]
If $|a_1|=1$, the maximum-modulus principle makes $h$ constant, and
the same inequality holds with $a_2=0$.

Since $f=(1+q)/(1-q)$,
$\log f=2a_1\zeta+2a_2\zeta^2+O(\zeta^3)$ near zero.
The second derivative of $\operatorname{artanh}$ is zero at zero, so
\[
 K\beta=2a_1,\qquad K^2\beta_x=4a_2.
\]
It follows that
$K^2(\beta^2+|\beta_x|)=4(a_1^2+|a_2|)\le4$, which proves
\eqref{eq:strip-log-budget}.
\end{proof}

\begin{theorem}[Carrier curvature and relative remainder jets]\label{thm:jets}
For every finite prime set containing $2$, put
$\alpha=(\log A_{\mathcal P})_x$. Then
\begin{align}
 |(\log A_{\mathcal P})_x|&\le\frac18,\label{eq:log-bound1}\\
 \alpha^2+|(\log A_{\mathcal P})_{xx}|&\le\frac1{64}.
       \label{eq:log-bound2}
\end{align}
In particular, $|(\log A_{\mathcal P})_{xx}|\le1/64$.
For a remainder starting at $m\ge2$, define
\begin{equation}\label{eq:remainder-slope}
 d_m=\frac18+\frac1{2m^2}.
\end{equation}
Its relative error $r_{\mathcal P}=R_{\mathcal P}/A_{\mathcal P}$
satisfies
\begin{align}
 |r_{\mathcal P,x}|&\le d_mr_{\mathcal P}\le d_m\rho_m,
       \label{eq:r1}\\
 |r_{\mathcal P,xx}|&\le d_m^2r_{\mathcal P}\le d_m^2\rho_m.
       \label{eq:r2}
\end{align}
For $\mathcal P=\{2,3\}$, one has $m=5$, $d_5=29/200$, and
\begin{equation}\label{eq:three-numbers}
 0<r_{\mathcal P}<7.3\times10^{-25},\quad
 |r_{\mathcal P,x}|<1.06\times10^{-25},\quad
 |r_{\mathcal P,xx}|<1.54\times10^{-26}.
\end{equation}
\end{theorem}
\begin{proof}
Apply Lemma~\ref{lem:strip-Schwarz} to $A_{\mathcal P}$ with $m=2$
to obtain \eqref{eq:log-bound1}--\eqref{eq:log-bound2}. Next put
$p=(\log R_{\mathcal P})_x$ and $q=(\log A_{\mathcal P})_x$.
The same lemma gives
\[
 p^2+|p_x|\le\frac1{4m^4},\qquad
 q^2+|q_x|\le\frac1{64}.
\]
Since $r_{\mathcal P}>0$, ordinary differentiation yields
\[
 \frac{r_{\mathcal P,x}}{r_{\mathcal P}}=p-q,\qquad
 \frac{r_{\mathcal P,xx}}{r_{\mathcal P}}
 =p_x-q_x+(p-q)^2.
\]
Thus $|r_{\mathcal P,x}|/r_{\mathcal P}\le|p|+|q|\le d_m$.
For the second derivative, keep the squared terms in both budgets:
\begin{align*}
 \frac{|r_{\mathcal P,xx}|}{r_{\mathcal P}}
 &\le |p_x|+|q_x|+p^2+q^2-2pq\\
 &\le\frac1{4m^4}+\frac1{64}-2pq
 \le\frac1{4m^4}+\frac1{64}+\frac1{8m^2}=d_m^2.
\end{align*}
Theorem~\ref{thm:relative} supplies $r_{\mathcal P}\le\rho_m$.
The displayed numerical inequalities follow from the scalar enclosure of
$\rho_5$ in Section~\ref{sec:certification}; even its rounded upper bound
$7.3\times10^{-25}$ gives
$d_5\rho_5<1.0585\times10^{-25}$ and
$d_5^2\rho_5<1.534825\times10^{-26}$.
\end{proof}

\subsection{A relative error throughout the unit strip}

The unit-strip lower bound also gives an explicit error estimate for
complex arguments. Define, using \eqref{eq:SQ} and \eqref{eq:kappa-strip},
\begin{equation}\label{eq:strip-relative-constants}
 I_m=\frac{S_m}{4\pi m^2-10},\qquad
 \rho_{s,m}=\frac{4097I_m+2Q_m+S_m}{\kappa_s}
 \qquad(m\ge2).
\end{equation}
The denominator defining $I_m$ is positive for $m\ge2$.

\begin{theorem}[Uniform complex relative approximation]\label{thm:strip-relative}
If $2\in\mathcal P$ and every index in $R_{\mathcal P}$ is at least
$m\ge2$, then
\begin{equation}\label{eq:strip-relative}
 \left|\frac{R_{\mathcal P}(t,z)}{A_{\mathcal P}(t,z)}\right|
 \le\rho_{s,m}
 \qquad(0\le t\le1/2,\ |\Im z|\le1).
\end{equation}
For $\mathcal P=\{2,3\}$,
\begin{equation}\label{eq:strip-number}
 \rho_{s,5}<8.5\times10^{-25}.
\end{equation}
\end{theorem}
\begin{proof}
Let $g(u)=e^{tu^2}\sum_{\substack{n>1\\(n,Q)=1}}\phi_n(u)$ and write
$R=R_{\mathcal P}$. The source majorant in
Lemma~\ref{lem:strip-tail}, with $Y=1$, gives
\begin{equation}\label{eq:strip-integral-mass}
 \int_0^\infty g(u)e^u\dd u\le I_m,\qquad |R(t,z)|\le I_m
 \quad(|\Im z|\le1).
\end{equation}
Set $q=-g'(0)$ and $s=g(0)$, so $q\le Q_m$, $s\le S_m$ and
$g''>0$ by Lemma~\ref{lem:mass} and the convexity proof. Two integrations
by parts give, for $z\ne0$,
\[
 z^2R(t,z)=q-\int_0^\infty g''(u)\cos(zu)\dd u.
\]
All boundary terms vanish at infinity, also after multiplication by
$e^u$. A second integration-by-parts calculation gives
\[
 \int_0^\infty g''(u)e^u\dd u
 =q+s+\int_0^\infty g(u)e^u\dd u.
\]
Since $|\cos(zu)|\le e^u$ for $|\Im z|\le1$, these identities imply
\begin{equation}\label{eq:strip-inverse-square}
 |z|^2|R(t,z)|\le2Q_m+S_m+I_m.
\end{equation}
The inequality is also true at $z=0$. Write $z=x+iy$ and use
$x^2\le|z|^2$. Equations \eqref{eq:strip-integral-mass} and
\eqref{eq:strip-inverse-square} yield
\[
 (4096+x^2)|R(t,z)|\le4097I_m+2Q_m+S_m.
\]
Finally $|A_{\mathcal P}|\ge\Re A_{\mathcal P}$, and
Proposition~\ref{prop:strip-lower} proves \eqref{eq:strip-relative}.
The strict numerical inequality \eqref{eq:strip-number} is certified
from the explicit constants in Section~\ref{sec:certification}.
\end{proof}

Both $H/A_{\mathcal P}$ and $R_{\mathcal P}/A_{\mathcal P}$ are
holomorphic on a neighborhood of the unit strip because the carrier has
positive real part in the larger strip $|\Im z|\le4\pi$. The absolute
lower bounds and displayed numerical tolerances in this section use the
original theta coefficient selectors. By contrast,
Theorem~\ref{thm:strip-positive} and Lemma~\ref{lem:strip-Schwarz}
hold for every bounded nonnegative higher-mode coefficient sequence.
They therefore give the same logarithmic-derivative budgets, and the
same factors $d_m,d_m^2$ in the relative derivative estimates, for
nonzero positive-coefficient carriers and remainders. These estimates
control an auxiliary quotient; they do not supply a signed margin for
the complete transform near a zero.

\begin{remark}[A larger auxiliary time range]\label{rem:time24}
All estimates in this section also hold for $0\le t\le24$. Indeed,
$2tu\le48u$ and
\[
 a\ge12(1+4u)>48u+9+\frac47
\]
still give $g'/g\le-3a$ in the convexity proof. The bound
$g''/g\ge7a^2$ uses only $t\ge0$ once this first-derivative estimate
is available. The local lower-bound proof is unchanged, and the source
majorant in Lemma~\ref{lem:strip-tail} remains valid because
$24<8\lambda_2\le8\lambda_n$ for $n\ge2$. The strip and logarithmic
derivative proofs therefore apply with the same constants. We retain
$0\le t\le1/2$ as the main range; this larger auxiliary range gives
no bound for the Newman constant.
\end{remark}
